\section{Litigation-Grade Recommendation}
\label{sec:litigation}

\subsection{The Two-Tier Standard}

Based on the 50-state table (Section~\ref{sec:table}), we recommend
a two-tier minimum ensemble standard for litigation-grade redistricting
analysis:

\begin{table}[h]
\centering
\caption{Litigation-grade minimum ensemble sizes by delegation size.
         Steps: total chain steps. ESS: expected effective sample size
         at the minimum step count. Power at 2pp: power to detect a
         2-percentage-point partisan advantage. Power at 10pp: power
         to detect a 10pp advantage.}
\label{tab:litigation-tiers}
\begin{tabular}{lccccc}
\toprule
Tier & Delegation & States & Min.\ steps & Min.\ ESS & Power at 10pp \\
\midrule
Tier 1 & $k \leq 17$ & 40 states & 20,000 & $\geq 1{,}000$ & $>99\%$ \\
Tier 2 & $k > 17$ & 7 states & 50,000 & $\geq 1{,}500$ & $>99\%$ \\
\bottomrule
\end{tabular}
\end{table}

The 40 Tier-1 states and 7 Tier-2 states are identified in Table~\ref{tab:state-table}.
The 3 non-redistricting states (AK, VT, WY at $k = 1$) do not apply.

\subsection{Rationale for 20,000 Steps (Tier 1)}

For NC ($k = 14$, the largest commonly litigated Tier-1 state):
\[
  \ess(20{,}000) = 20{,}000 \cdot \frac{1 - 0.87}{1 + 0.87} \approx 1{,}390.
\]
This provides:
\begin{itemize}
\item Convergence: well beyond $\rhat < 1.1$ (achieved at $\approx 2{,}000$ steps)
\item Tail precision: 95\% CI width on a 1st-percentile claim is
  $\pm 0.53$ percentage points (vs.\ $\pm 0.75$ pp at 10,000 steps)
\item Power: $>99\%$ for any practically significant gerrymandering effect
\end{itemize}

For WI ($k = 8$):
\[
  \ess(20{,}000) \approx 20{,}000 \cdot \frac{0.15}{1.85} \approx 1{,}622.
\]
Even stronger guarantees.

\subsection{Rationale for 50,000 Steps (Tier 2)}

For CA ($k = 52$):
\[
  \ess(50{,}000) = 50{,}000 \cdot \frac{1 - 0.94}{1 + 0.94} \approx 1{,}546.
\]
G.4 certifies convergence ($\rhat < 1.1$) at 50,000 steps for CA.
ESS $\approx 1{,}546$ provides power $>99\%$ for any 10pp effect.

For TX ($k = 38$):
\[
  \ess(50{,}000) = 50{,}000 \cdot \frac{1 - 0.92}{1 + 0.92} \approx 2{,}083.
\]
G.4 certifies convergence at 25,000 steps for TX; 50,000 steps provides
a $2\times$ safety margin.

\subsection{Feasibility}

At \bisect{}'s H.2 ensemble module throughput of 50,000 steps per second:
\begin{itemize}
\item 20,000 steps: $\leq 0.4$ seconds per state
\item 50,000 steps: $\leq 1.0$ second per state
\item Full 50-state Tier-1 run: $\leq 20$ seconds
\item Full 50-state Tier-2 run (large states only): $\leq 7$ seconds
\end{itemize}

The recommended minimum ensemble sizes are not computationally burdensome
with current redistricting software.
The constraint is not computation time but the awareness that power
analysis is required before reporting percentile claims.

\subsection{Required Disclosures in Expert Reports}

An expert using a redistricting ensemble for litigation purposes should
disclose the following:
\begin{enumerate}
\item Chain method (\recom{}, Forest-ReComm, SMC, etc.)
\item Number of steps (nominal $n$) and chains (parallel chains count)
\item Autocorrelation diagnostic: $\rho_1$ and $\ess = n^*$
\item Convergence certification: $\rhat$ for the primary test statistic
\item Power: the power at the claimed effect size given the certified $\ess$
\item Significance level $\alpha$ and whether multiple testing correction was applied
\end{enumerate}

Without these disclosures, a court cannot assess whether the ensemble
evidence is reliable.
S.4~\citep{s4paper} extends this list to the full test battery.
